\appendix
\section{Código Computacional}
\label{anexo:codigo_python}

O código a seguir apresenta a implementação completa do script em Python, orientado a objetos, utilizado para o cálculo das expressões intermediárias das grandezas cinemáticas e para a geração dos triângulos de velocidades apresentados na metodologia do trabalho.

% Configuração para aceitar acentos em português dentro do código Python
\lstset{literate=
  {á}{{\'a}}1 {é}{{\'e}}1 {í}{{\'i}}1 {ó}{{\'o}}1 {ú}{{\'u}}1
  {Á}{{\'A}}1 {É}{{\'E}}1 {Í}{{\'I}}1 {Ó}{{\'O}}1 {Ú}{{\'U}}1
  {à}{{\`a}}1 {è}{{\`e}}1 {ì}{{\`i}}1 {ò}{{\`o}}1 {ù}{{\`u}}1
  {À}{{\`A}}1 {È}{{\'E}}1 {Ì}{{\`I}}1 {Ò}{{\`O}}1 {Ù}{{\`U}}1
  {ä}{{\"a}}1 {ë}{{\"e}}1 {ï}{{\"i}}1 {ö}{{\"o}}1 {ü}{{\"u}}1
  {Ä}{{\"A}}1 {Ë}{{\"E}}1 {Ï}{{\"I}}1 {Ö}{{\"O}}1 {Ü}{{\"U}}1
  {â}{{\^a}}1 {ê}{{\^e}}1 {î}{{\^i}}1 {ô}{{\^o}}1 {û}{{\^u}}1
  {Â}{{\^A}}1 {Ê}{{\^E}}1 {Î}{{\^I}}1 {Ô}{{\^O}}1 {Û}{{\^U}}1
  {ã}{{\~a}}1 {ẽ}{{\~e}}1 {ĩ}{{\~i}}1 {õ}{{\~o}}1 {ũ}{{\~u}}1
  {Ã}{{\~A}}1 {Ẽ}{{\~E}}1 {Ĩ}{{\~I}}1 {Õ}{{\~O}}1 {Ũ}{{\~U}}1
  {œ}{{\oe}}1 {Œ}{{\OE}}1 {æ}{{\ae}}1 {Æ}{{\AE}}1 {ß}{{\ss}}1
  {ç}{{\c c}}1 {Ç}{{\c C}}1 {ø}{{\o}}1 {å}{{\r a}}1 {Å}{{\r A}}1
}

\begin{lstlisting}[language=Python, caption={Script de simulação cinemática da Lewis VL750}]
import pandas as pd
import numpy as np
import matplotlib.pyplot as plt

class TurbomachineAnalyzer:
    """
    Representa um analisador de geometria e desempenho de turbomáquinas.

    Esta classe fornece métodos para calcular e recuperar propriedades geométricas,
    componentes de velocidade e características de escoamento para seções específicas de uma
    turbomáquina. É destinada à análise da dinâmica de fluidos e aos parâmetros de projeto
    de bombas, turbinas e máquinas similares.

    Atributos
    ---------
    D : float
        Diâmetro externo da turbomáquina.
    d : float
        Diâmetro interno da turbomáquina.
    N : float
        Velocidade de rotação em rotações por minuto (RPM).
    Q : float
        Vazão volumétrica através da turbomáquina.
    alpha_4 : float
        Ângulo de incidência do escoamento (velocidade absoluta) na seção 4, em graus.
    beta_5 : float
        Ângulo de saída do escoamento (velocidade relativa) na seção 5, em graus.

    Métodos
    -------
    __init__(D, d, N, Q, alpha_4, beta_5)
        Inicializa o analisador com os parâmetros geométricos e operacionais fornecidos.

    get_geometry()
        Recupera e retorna as propriedades geométricas e operacionais como um dicionário.

    get_section4()
        Recupera os dados da seção 4 como um dicionário.

    get_section5()
        Recupera os dados da seção 5 como um dicionário.

    get_energy_transfer()
        Calcula a transferência de energia (Teorema de Euler).

    plot_velocity_triangles(show=True)
        Gera e plota os triângulos de velocidades da sucção e do recalque.
    """

    def __init__(self, D, d, N, Q, alpha_4, beta_5):
        self.D = float(D)
        self.d = float(d)
        self.N = float(N)
        self.Q = float(Q)
        self.alpha_4 = float(alpha_4)
        self.beta_5 = float(beta_5)

        self._compute_geometry()
        self._compute_section4()
        self._compute_section5()

    def _compute_geometry(self):
        """
        Calcula diversas propriedades geométricas e de desempenho.

        Este método calcula atributos geométricos como raio, raio médio, diâmetro
        e outros usando as dimensões de entrada e parâmetros operacionais do objeto.
        É essencial para determinar as características derivadas do sistema.

        Atributos calculados
        --------------------
        R : float
            Raio externo calculado como metade do diâmetro externo (D).
        r : float
            Raio interno calculado como metade do diâmetro interno (d).
        rm : float
            Raio médio calculado como a média geométrica entre R e r.
        dm : float
            Diâmetro médio calculado como o dobro do raio médio.
        U : float
            Velocidade periférica (tangencial) calculada usando o diâmetro médio
            e a velocidade de rotação (N).
        A : float
            Área da seção transversal determinada pelos diâmetros externo e interno.
        Cm : float
            Velocidade meridiana média calculada como a razão entre a vazão
            volumétrica (Q) e a área da seção transversal (A).
        """
        self.R = self.D / 2.0
        self.r = self.d / 2.0
        self.rm = (self.R + self.r) / 2.0
        self.dm = 2.0 * self.rm
        self.U = np.pi * self.dm * self.N / 60.0
        self.A = (np.pi / 4.0) * (self.D**2 - self.d**2)
        self.Cm = self.Q / self.A

    def _compute_section4(self):
        """
        Realiza os cálculos para os triângulos de velocidades e propriedades da seção de entrada (Estação 4).

        Este método inclui cálculos relevantes para o escoamento, como as componentes
        tangencial e meridiana da velocidade, velocidade absoluta, velocidade relativa e seus
        respectivos ângulos.

        Atributos calculados
        --------------------
        Cm4 : float
            Componente meridiana da velocidade na seção 4.
        Cu4 : float
            Componente tangencial da velocidade absoluta na seção 4. Calculada com base em `alpha_4`.
        C4 : float
            Magnitude (módulo) da velocidade absoluta na seção 4.
        Wu4 : float
            Componente tangencial da velocidade relativa à pá na seção 4.
        Wm4 : float
            Componente meridiana da velocidade relativa à pá na seção 4.
        W4 : float
            Magnitude (módulo) da velocidade relativa na seção 4.
        beta_4_calc : float
            Ângulo da velocidade relativa `W4` em relação à direção meridiana, em graus.
        """
        self.Cm4 = self.Cm

        # Se alpha for 90, entra reto (Cu = 0). Senão, calcula a componente tangencial.
        if np.isclose(self.alpha_4, 90.0):
            self.Cu4 = 0.0
        else:
            self.Cu4 = self.Cm4 / np.tan(np.radians(self.alpha_4))

        self.C4 = np.sqrt(self.Cm4**2 + self.Cu4**2)

        # Geometria do triângulo (W fecha C com U)
        self.Wu4 = self.U - self.Cu4
        self.Wm4 = self.Cm4
        self.W4 = np.sqrt(self.Wu4**2 + self.Wm4**2)

        self.beta_4_calc = np.degrees(np.arctan2(self.Wm4, self.Wu4))
\end{lstlisting}